\section*{第 \ref{ch:b3:product} 章——乘积测度，
富比尼、变量代换}

\begin{solution}{exo:b3:product:1}
$\Delta$ 在 $\intcc01^2$ 中是闭集，因而是博雷尔集，且 $\mathcal
B(\intcc01^2)$ 是乘积 $\sigma$-代数
（\cref{prop:b3:product:sections}(b)）。按一种次序累次积分：
\[
\int_{\intcc01}\Bigl(\int \mathbf 1_\Delta(x,y)
\,\dd\lambda(y)\Bigr)\dd\mu(x)
= \int \lambda(\{x\})\,\dd\mu(x) = 0 ;
\]
另一种次序：
\[
\int_{\intcc01}\Bigl(\int\mathbf
1_\Delta(x,y)\,\dd\mu(x)\Bigr)\dd\lambda(y)
= \int \mu(\{y\})\,\dd\lambda(y) = \int 1\,\dd\lambda = 1 .
\]
失效的假设是 $\sigma$-有限性：计数\omterm{def:b3:measure:measure}{测度} $\mu$ 在
不可数集 $\intcc01$ 上不满足它。
没有任何由 $\mu$ 有限的集合构成的可数族能覆盖它。
\end{solution}

\begin{solution}{exo:b3:product:2}
在 $\intcc0A\times\intoo0{+\infty}$ 上：
$\int_0^A\int_0^\infty\eu^{-xy}\abs{\sin x}\,\dd y\,\dd x =
\int_0^A\frac{\abs{\sin x}}x\dd x \leq A < \infty$（托内利定理，对绝对值）
富比尼定理可用，且由于
$\int_0^\infty\eu^{-xy}\dd y = \frac1x$，于是
\[
\int_0^A\frac{\sin x}x\dd x
= \int_0^\infty\Bigl(\int_0^A\eu^{-xy}\sin x\,\dd
x\Bigr)\dd y
= \int_0^\infty \frac{1 - \eu^{-Ay}(\cos A + y\sin A)}{1 +
y^2}\,\dd y
\]
（内层积分：$\operatorname{Im}\int_0^A\eu^{(\iu -
y)x}\dd x$，直接算出）。当 $A \to \infty$ 时，
修正项的界为 $\int_0^\infty\eu^{-Ay}\frac{1 +
y}{1 + y^2}\dd y \leq \frac32\int_0^\infty\eu^{-Ay}\dd y =
\frac3{2A} \to 0$；主项为
$\int_0^\infty\frac{\dd y}{1+y^2} = \frac\pi2$。于是
$\int_0^\infty\frac{\sin x}x\dd x = \frac\pi2$——由富比尼定理得出的
狄利克雷积分。
\end{solution}

\begin{solution}{exo:b3:product:3}
(a) 层饼公式（\cref{prop:b3:product:layercake}）：$\int
f\,\dd\mu = \int_0^\infty\mu(f > t)\,\dd t$，且 $t \mapsto
\mu(f > t)$ 为非增函数。在 $[n-1, n]$ 上：$\mu(f \geq n)
\leq \mu(f > t) \leq \mu(f > n - 1) \leq \mu(f \geq n - 1)$；
对各单位区间上的积分求和：
\[
\sum_{n\geq1}\mu(f \geq n) \leq \int f\,\dd\mu \leq
\sum_{n\geq1}\mu(f \geq n - 1) = \mu(f \geq 0) +
\sum_{n\geq1}\mu(f\geq n) \leq \mu(X) + \sum_{n\geq1}\mu(f\geq
n).
\]

(b) 对 $\abs f$ 应用 (a)：积分的有限性与
级数的有限性等价（多出的 $\mu(X)$ 有限）。
\end{solution}

\begin{solution}{exo:b3:product:4}
由于 $f(x,y) = \partial_y\bigl(\frac{y}{x^2+y^2}\bigr)$ 在
$x \ne 0$ 时成立：
\[
\int_0^1 f(x, y)\,\dd y = \frac{1}{x^2 + 1}
\ \Longrightarrow\
\int_0^1\!\!\int_0^1 f\,\dd y\,\dd x =
\int_0^1\frac{\dd x}{1 + x^2} = \frac\pi4 ;
\]
由反对称性 $f(y,x) = -f(x,y)$，另一种次序给出
$-\frac\pi4$。取绝对值，对 $0 < y < x$，
\[
\int_0^x f(x,y)\,\dd y = \Bigl[\frac{y}{x^2 +
y^2}\Bigr]_0^x = \frac1{2x},
\quad\text{且 } f \geq 0 \text{ 在该处，因此}\quad
\int_0^1\!\!\int_0^1\abs f \geq \int_0^1\frac{\dd x}{2x} =
+\infty .
\]
与富比尼定理并不矛盾：其假设 $f \in L^1$
不成立，两个累次积分恰好是两个不同的
数值。
\end{solution}

\begin{solution}{exo:b3:product:5}
(a) $(\mathbf 1_{\intcc01}*\mathbf 1_{\intcc01})(x) =
\lambda\bigl(\intcc01\cap\intcc{x-1}x\bigr)$：$0$，当 $x
\notin \intcc02$ 时；$x$，当 $0 \leq x \leq 1$ 时；$2 - x$，当 $1
\leq x \leq 2$：帐篷函数。再卷积一次，得到 $\mathcal
C^1$ 的分段二次凸峰，位于 $\intcc03$ 上（二次
B 样条）：每次卷积使光滑度提高一级——
光滑化原理，正是 \cref{ch:b3:lp} 中磨光子的依据。

(b) 若 $x \notin \overline{\operatorname{supp}f +
\operatorname{supp}g}$，则存在围绕 $x$ 的一个球，
与和集不相交；对 $y \in \operatorname{supp}g$，$x - y
\notin\operatorname{supp}f$，故被积函数
恒等于零：$f * g = 0$ 在 $x$ 附近。

(c) 对 $x_n \to x$：$(f*g)(x_n) = \int
f(y)g(x_n - y)\,\dd y$；被积函数逐点收敛
（由 $g$ 的连续性），并被 $\norm
g_\infty\,\abs f \in L^1$ 控制：由控制收敛定理，$(f*g)(x_n) \to (f*g)(x)$。
\end{solution}

\begin{solution}{exo:b3:product:6}
(a) 由富比尼定理与归纳法，沿最后一个
坐标切片：
\[
\lambda_d(\Delta_d) = \int_0^1
\lambda_{d-1}\bigl((1 - t)\,\Delta_{d-1}\bigr)\,\dd t
= \lambda_{d-1}(\Delta_{d-1})\int_0^1(1 - t)^{d-1}\dd t
= \frac{\lambda_{d-1}(\Delta_{d-1})}{d},
\]
利用伸缩性质 $\lambda_{d-1}(\rho A) =
\rho^{d-1}\lambda_{d-1}(A)$
（见\cref{thm:b3:product:linearchange}）；
取 $\lambda_1(\Delta_1) = 1$，得体积 $\frac1{d!}$。

(b) $v_2 = \pi/\Gamma(2) = \pi$；$v_3 =
\pi^{3/2}/\Gamma(\frac52) = \pi^{3/2}/(\frac32\cdot\frac12
\sqrt\pi) = \frac{4\pi}3$。椭球为 $T(B(0,1))$，其中
$T = \operatorname{diag}(a_1, \dots, a_d)$：
由\cref{thm:b3:product:linearchange}可知体积为
$v_d\prod a_i$。
\end{solution}

\begin{solution}{exo:b3:product:7}
在 $\R^2$ 中，取\omterm{ex:b3:product:polar}{极坐标}（见\cref{ex:b3:product:polar}）：
\[
\int_{B(0,1)}\frac{\dd x\dd y}{(x^2+y^2)^s}
= 2\pi\int_0^1 r^{1 - 2s}\,\dd r,
\qquad
\int_{\R^2\setminus B(0,1)} = 2\pi\int_1^\infty r^{1-2s}\dd r:
\]
有限当且仅当 $1 - 2s > -1$（$s < 1$），外部积分则要求 $1 - 2s < -1$（$s >
1$）。在 $\R^d$ 中避免使用球坐标，改用层饼
表示：$\lambda_d(\{\norm x^{-2s} > t\}\cap B(0,1)) =
\lambda_d(B(0, \min(1, t^{-1/2s}))) = v_d\min(1, t^{-d/2s})$，
且 $\int_0^\infty v_d\min(1, t^{-d/(2s)})\dd t < \infty$ 当且仅当
$\frac d{2s} > 1$，即 $s < \frac d2$；外部积分
收敛当且仅当 $s > \frac d2$（在
补集区域上作同样的计算）。
\end{solution}

\begin{solution}{exo:b3:product:8}
由托内利定理（被积函数非负）及变量替换
$(x, y) = \Phi(u, v) = (uv,\ u(1-v))$，它是 $\mathcal C^1$
微分\omterm{def:b3:topology:continuity}{同胚}，将 $\intoo0\infty\times\intoo01$ 映到开象限上，
并有
\[
\det D\Phi = \det\begin{pmatrix} v & u\\ 1 - v & -u
\end{pmatrix} = -uv - u(1 - v) = -u,
\qquad \abs{\det} = u :
\]
\[
\Gamma(p)\Gamma(q) =
\iint x^{p-1}y^{q-1}\eu^{-x-y}\dd x\,\dd y
= \iint (uv)^{p-1}\bigl(u(1{-}v)\bigr)^{q-1}\eu^{-u}\,u\,
\dd u\,\dd v
= \Gamma(p + q)\,B(p, q).
\]
将 $t = \sin^2\theta$ 代入 $B(p,q)$，得
$2\int_0^{\pi/2}\sin^{2p-1}\theta\cos^{2q-1}\theta\,\dd\theta =
B(p, q)$。沃利斯积分：$W_n =
\int_0^{\pi/2}\sin^n\theta\,\dd\theta = \frac12B\bigl(\frac{n +
1}2, \frac12\bigr) = \frac{\Gamma(\frac{n+1}2)\sqrt\pi}
{2\,\Gamma(\frac n2 + 1)}$——例如 $W_{2n} =
\frac\pi2\cdot\frac{(2n)!}{4^n(n!)^2}$，利用
$\Gamma(n + \frac12) = \frac{(2n)!}{4^nn!}\sqrt\pi$。
\end{solution}

\begin{solution}{exo:b3:product:9}
示性函数：$\int\mathbf 1_B\,\dd(T_*\mu) = T_*\mu(B) =
\mu(T^{-1}B) = \int\mathbf 1_B\circ T\,\dd\mu$；由线性
可推至\omterm{def:b3:lebesgue:simple}{简单函数} $g$，再由单调收敛定理推至 $g \geq 0$——转移
公式。它纯属形式：它把关于
$T_*\mu$ 的积分重新表述，但对 $T_*\mu$ 究竟\emph{是什么}并未涉及。真正的
内容在于 \cref{thm:b3:product:linearchange} 与
\cref{thm:b3:product:changeofvar} 所给出的如下恒等式：
\[
\Phi_*\bigl(\lambda_d\restriction_U\bigr) =
\abs{\det D\Phi^{-1}}\,\lambda_d\restriction_V
\quad\text{（一个密度测度）},
\]
即计算\omterm{def:b3:measure:lebesgueouter}{勒贝格测度}的\omterm{thm:b3:product:ballvolume}{像测度}——
其解析输入是 $\Phi$ 的微分几何，
而非\omterm{def:b3:measure:measure}{测度}论的形式主义。
\end{solution}

\begin{solution}{exo:b3:product:10}
由托内利定理，高斯被积函数可以分解，取 $G_1 =
\int_\R\eu^{-s^2}\dd s = \sqrt\pi$ 及 $\int_\R
s^2\eu^{-s^2}\dd s = \frac{\sqrt\pi}2$（分部积分）：
\[
\int_{\R^d}x_1^2\,\eu^{-\norm x^2}\dd x =
\frac{\sqrt\pi}2\;\pi^{(d-1)/2} = \frac{\pi^{d/2}}2,
\qquad
\int_{\R^d}\norm x^2\eu^{-\norm x^2}\dd x =
d\cdot\frac{\pi^{d/2}}2
\]
（由对称性，$\norm x^2 = \sum_ix_i^2$ 的 $d$ 个项贡献相等，
即一致性检验）。对于归一化\omterm{def:b3:measure:measure}{测度}
$\pi^{-d/2}\eu^{-\norm x^2}\dd x$，其二阶矩为
$\frac d2$。
\end{solution}

\begin{solution}{exo:b3:product:11}
(a) 记 $H = \Phi^{-1}(\intoo0\infty)$，其中 $\Phi(x, y) = f(x) -
y$ 与 $\{y > 0\}$ 的交集可测，因为 $(x, y)
\mapsto f(x)$ 与 $(x,y)\mapsto y$ 在乘积上都可测
（它们均为与投影的复合）。
$x$-截面（对 $H$）为 $\intoo0{f(x)}$，\omterm{def:b3:measure:measure}{测度}为 $f(x)$：
由托内利定理对截面积分，
\[
\lambda_{d+1}(H) = \int_{\R^d}\lambda_1\bigl(\intoo0{f(x)}
\bigr)\,\dd x = \int_{\R^d}f\,\dd\lambda_d .
\]

(b) 图像为 $\{(x,y) : y \geq f(x)\} \cap \{y \leq
f(x)\}$，可测；其 $x$-截面为单点集，
\omterm{def:b3:measure:measure}{测度}为 $0$，由托内利定理得 $\lambda_{d+1}(\text{图像}) =
\int 0 = 0$。

(c) $S^{d-1}$ 是两图像 $y =
\pm\sqrt{1 - \abs{x'}^2}$ 的并，定义在 $\R^{d-1}$ 的单位球上，
（沿最后一个坐标分裂）：两个零集的并，由
(b)，仍为零集。
\end{solution}

\begin{solution}{exo:b3:product:12}
在 $\intoo01^2$ 上，$\frac1{1 - xy} = \sum_{n\geq0}(xy)^n$，
各项非负，由托内利定理可逐项积分，
\[
\iint\frac{\dd x\,\dd y}{1 - xy}
= \sum_{n\geq0}\Bigl(\int_0^1x^n\dd x\Bigr)
\Bigl(\int_0^1y^n\dd y\Bigr)
= \sum_{n\geq0}\frac1{(n+1)^2} = \zeta(2) .
\]
对交错情形，$\frac1{1 + xy} =
\sum_n(-1)^n(xy)^n$ 不是正项级数；但
\emph{绝对值}级数的积分为 $\zeta(2) <
\infty$，故富比尼定理（此时\omterm{def:b3:lebesgue:l1}{可积}性已确立）适用：
$\iint\frac{\dd x\dd y}{1 + xy} =
\sum_n\frac{(-1)^n}{(n+1)^2}$。级数恒等式为：
\[
\sum_{n\geq1}\frac{(-1)^{n-1}}{n^2}
= \sum_{n\geq1}\frac1{n^2} - 2\sum_{k\geq1}\frac1{(2k)^2}
= \zeta(2) - \frac{\zeta(2)}2 = \frac{\zeta(2)}2 .
\]
利用 $\zeta(2) = \frac{\pi^2}6$（见\cref{pb:b3:spectral:1}）：
两积分分别为 $\frac{\pi^2}6$ 与 $\frac{\pi^2}{12}$。
在第一个计算中，托内利定理的正性是整个论证的关键
——无需在交换之前检验\omterm{def:b3:lebesgue:l1}{可积}性；
在第二个计算中，绝对值级数的正性
正是\emph{保证}\omterm{def:b3:lebesgue:l1}{可积}性的依据，使富比尼定理
得以应用于带符号的级数。
\end{solution}

\begin{solution}{pb:b3:product:1}
\textbf{1.} 令 $x = t + \sqrt t\,u$（$\dd x = \sqrt
t\,\dd u$；当 $x$ 取遍 $\intoo0\infty$ 时，$u$ 取遍
$\intoo{-\sqrt t}\infty$）：
\[
\Gamma(t{+}1) = \int_0^\infty x^t\eu^{-x}\dd x
= t^t\eu^{-t}\sqrt t\int_{-\sqrt t}^{\infty}
\exp\Bigl(t\ln\Bigl(1 + \frac u{\sqrt t}\Bigr) - \sqrt
t\,u\Bigr)\dd u,
\]
由 $x^t = t^t\exp\bigl(t\ln(1 + u/\sqrt t)\bigr)$ 与
$\eu^{-x} = \eu^{-t}\eu^{-\sqrt tu}$ 即得上式。

\textbf{2.} 固定 $u$，$t \to \infty$ 时：$t\ln(1 +
u/\sqrt t) - \sqrt tu = t\bigl(\frac u{\sqrt t} -
\frac{u^2}{2t} + o(\frac1t)\bigr) - \sqrt tu = -\frac{u^2}2 +
o(1)$：$g_t(u) \to \eu^{-u^2/2}$。

\textbf{3.} 设 $\psi_1(h) = \varphi(h) + \frac{h^2}4$，
定义于 $\intoc{-1}1$ 上，$\psi_1(0) = 0$，$\psi_1'(h) = \frac1{1+h} -
1 + \frac h2 = \frac{h(h-1)}{2(1+h)}$，此式 $\geq 0$ 在
$\intoc{-1}0$ 上成立，$\leq 0$ 在 $\intcc01$ 上成立：$\psi_1 \leq 0$，
即 $\varphi(h) \leq -h^2/4$ 在此区间成立。令 $\psi_2(h) =
\varphi(h) + ch$，定义于 $\intco1\infty$ 上，$c = 1 - \ln2$：$\psi_2(1)
= \ln2 - 1 + c = 0$，且 $\psi_2'(h) = c - \frac h{1+h} \leq c -
\frac12 < 0$，于是 $\varphi(h) \leq -ch$ 对 $h \geq 1$ 成立。现取 $t
\geq 1$，若 $\abs u \leq \sqrt t$，则 $g_t(u) =
\eu^{t\varphi(u/\sqrt t)} \leq \eu^{-t(u/\sqrt t)^2/4} =
\eu^{-u^2/4}$；若 $u \geq \sqrt t$，则 $t\,\varphi(u/\sqrt t) \leq
-ct\cdot\frac u{\sqrt t} = -c\sqrt t\,u \leq -cu$（因 $t \geq
1$），故 $g_t(u) \leq \eu^{-cu}$。于是 $g_t \leq \eu^{-u^2/4} +
\eu^{-cu}\mathbf 1_{u>0}$，\omterm{def:b3:lebesgue:l1}{可积}，且与 $t \geq
1$ 无关。

\textbf{4.} 由控制收敛定理，$\int_\R g_t(u)\dd u \to
\int_\R\eu^{-u^2/2}\dd u = \sqrt{2\pi}$
（利用\cref{ex:b3:product:polar}，再作
尺度变换 $u\mapsto\sqrt2\,u$）。结合第 1 问：
\[
\Gamma(t + 1) \sim \sqrt{2\pi t}\;\Bigl(\frac
t\eu\Bigr)^t,\qquad
n! \sim \sqrt{2\pi n}\,\Bigl(\frac n\eu\Bigr)^n .
\]

\textbf{5.} 由\cref{exo:b3:product:8}，$W_{2n} =
\frac\pi2\,\frac{(2n)!}{4^n(n!)^2} =
\frac\pi2\,4^{-n}\binom{2n}n$。由斯特林公式：
\[
\binom{2n}{n} = \frac{(2n)!}{(n!)^2}
\sim \frac{\sqrt{4\pi n}\,(2n/\eu)^{2n}}
{2\pi n\,(n/\eu)^{2n}} = \frac{4^n}{\sqrt{\pi n}} .
\]
于是 $W_{2n} \sim \frac12\sqrt{\pi/n}$，与
递推式 $W_n = \frac{n-1}nW_{n-2}$（该式蕴含 $W_n \sim
W_{n-2}$，并结合 $W_nW_{n-1}\cdot n = \frac\pi2$——
经典沃利斯关系——给出 $W_n \sim
\sqrt{\pi/(2n)}$；两种渐近式相容）。

\textbf{6.} $\Gamma(\frac d2 + 1) \sim \sqrt{2\pi\frac
d2}\,(\frac d{2\eu})^{d/2}$，于是
\[
v_d = \frac{\pi^{d/2}}{\Gamma(\frac d2 + 1)}
\sim \frac{1}{\sqrt{\pi d}}\Bigl(\frac{2\pi\eu}
d\Bigr)^{d/2} \longrightarrow 0
\]
比几何级数更快地趋于零（当 $d > 2\pi\eu \approx 17$ 时，每个因子
$< 1$ 且逐步减小）。数值上 $v_1 = 2$，$v_2 \approx
3.14$，$v_3 \approx 4.19$，$v_4 \approx 4.93$，$v_5 \approx
5.26$，$v_6 \approx 5.17$：最大值在 $d = 5$ 处取得。

\textbf{7.} 取 $k = \frac n2 + \frac{s\sqrt n}2$（整数，
$n$ 为偶数，$s$ 固定）：对 $2^{-n}\binom nk =
2^{-n}\frac{n!}{k!(n-k)!}$ 取对数，并对三个阶乘
应用斯特林公式。记 $k = \frac n2(1 + \varepsilon)$，$n - k =
\frac n2(1 - \varepsilon)$，其中 $\varepsilon = s/\sqrt n$：
\[
\ln\Bigl(2^{-n}\binom nk\Bigr)
= -\frac n2\bigl[(1{+}\varepsilon)\ln(1{+}\varepsilon) +
(1{-}\varepsilon)\ln(1{-}\varepsilon)\bigr]
+ \frac12\ln\frac{2}{\pi n(1 - \varepsilon^2)} + o(1),
\]
方括号内的量为 $\varepsilon^2 + O(\varepsilon^4) =
\frac{s^2}n + O(n^{-2})$：上式趋于 $-\frac{s^2}2 +
\frac12\ln\frac2{\pi n}$ 误差为 $o(1)$，即
\[
2^{-n}\binom nk \sim \sqrt{\frac{2}{\pi
n}}\;\eu^{-s^2/2} :
\]
抛掷硬币所呈现的高斯分布形态，已得到定量刻画——
棣莫弗–拉普拉斯的局部形式，有待在
\cref{ch:b3:clt}中加以整体化。

\textbf{8.} (i) 控制收敛定理把第 2 题中的逐点极限
化为积分的收敛，所用控制函数取自第 3 题。
(ii) 高斯积分计算出极限
$\int\eu^{-u^2/2} = \sqrt{2\pi}$——斯特林常数
$\sqrt{2\pi}$ \emph{正是}高斯积分。
(iii) 代换 $x = t + \sqrt tu$ 是一个仿射变量变换，
其依据是平移不变性以及\omterm{def:b3:measure:lebesgueouter}{勒贝格测度}的伸缩法则
（\cref{thm:b3:product:linearchange}，
取维数 $1$ 的情形）。

\textbf{9.} 展开两项：
\[
d_n - d_{n+1} = \ln\frac{n!}{(n+1)!} + \Bigl(n +
\frac32\Bigr)\ln(n+1) - \Bigl(n + \frac12\Bigr)\ln n - 1
= \Bigl(n + \frac12\Bigr)\ln\frac{n+1}n - 1,
\]
两项 $-\ln(n+1)$ 与 $(n + \frac32)\ln(n+1)$ 合并成
$(n + \frac12)\ln(n+1)$。

\textbf{10.} 取 $t = \frac1{2n+1}$，则 $\frac{1+t}{1-t} =
\frac{2n+2}{2n} = \frac{n+1}n$，且 $n + \frac12 =
\frac1{2t}$；奇幂级数 $\ln\frac{1+t}{1-t} =
2\sum_{k\geq0}\frac{t^{2k+1}}{2k+1}$ 给出
\[
\Bigl(n + \frac12\Bigr)\ln\frac{n+1}n
= \sum_{k\geq0}\frac{t^{2k}}{2k+1}
= 1 + \frac{t^2}3 + \frac{t^4}5 + \cdots
\]
减去 $1$。下界：只取第一项，
$\frac{t^2}3 = \frac1{3(2n+1)^2}$。上界：将所有分母降为
$3$，并对等比级数求和：
$\frac{t^2}{3(1 - t^2)} = \frac1{3((2n+1)^2 - 1)} =
\frac1{12n(n+1)} = \frac1{12n} - \frac1{12(n+1)}$.

\textbf{11.} 将上界从 $n$ 到 $\infty$ 求和
（用到 $d_m \to 0$）得 $d_n < \frac1{12n}$。下界：
$\frac1{12m+1} - \frac1{12(m+1)+1} =
\frac{12}{(12m+1)(12m+13)}$，且
\begin{align*}
\frac1{3(2m+1)^2} > \frac{12}{(12m+1)(12m+13)}
&\iff (12m+1)(12m+13) > 36(2m+1)^2 \\
&\iff 168m + 13 > 144m + 36,
\end{align*}
对 $m \geq 1$ 成立。对此可裂项相消的下界求和：
$d_n > \frac1{12n+1}$。取指数后得到经典的
$n!$ 夹逼估计。

\textbf{12.} (a) 相对误差 $= \eu^{d_n} - 1 <
\eu^{1/(12n)} - 1 < \frac{1.1}{12n}$，当 $n$ 充分大时；$<
10^{-6}$ 只要 $12n \geq 1.1\cdot10^6$，而所给的
$n \geq 83\,334$ 即已足够（$\frac1{12n} \leq 10^{-6}$
已蕴含它）。(b) $\log_{10}(100!) =
\frac12\log_{10}(200\pi) + 200 - 100\log_{10}\eu +
d_{100}\log_{10}\eu = 1.39906 + 200 - 43.42945 + 0.00036
\approx 157.96997$，故 $100! \approx 10^{0.96997} \cdot
10^{157} = 9.333\cdot10^{157}$；可保证的区间
$(\eu^{1/1201}, \eu^{1/1200})$ 的相对宽度
小于 $10^{-6}$——所谓“渐近”公式，在 $n =
100$ 时却是一件精密的工具。

\textbf{13.} 写 $\sin^n = \sin^{n-2} - \sin^{n-2}\cos^2$，
并对第二项分部积分（取 $u = \cos\theta$，
$\dd v = \sin^{n-2}\cos\theta\,\dd\theta$，$v =
\frac{\sin^{n-1}}{n-1}$）：
\[
\int_0^{\pi/2}\sin^{n-2}\cos^2 =
\Bigl[\cos\theta\,\frac{\sin^{n-1}\theta}{n-1}\Bigr]_0^{\pi/2}
+ \frac1{n-1}\int_0^{\pi/2}\sin^n = \frac{W_n}{n-1} .
\]
于是 $W_n = W_{n-2} - \frac{W_n}{n-1}$，即 $W_n =
\frac{n-1}nW_{n-2}$。由 $W_0 = \frac\pi2$，$W_1 = 1$ 可得：
\[
W_{2n} = \frac{(2n-1)!!}{(2n)!!}\cdot\frac\pi2
= \frac\pi2\binom{2n}n4^{-n},
\qquad
W_{2n+1} = \frac{(2n)!!}{(2n+1)!!}
= \frac{4^n}{(2n+1)\binom{2n}n},
\]
双阶乘的换算依据 $(2n)!! = 2^nn!$ 及
$(2n-1)!! = \frac{(2n)!}{2^nn!}$。最后 $nW_nW_{n-1} =
(n-1)W_{n-1}W_{n-2}$ 由递推式知其为常数，且等于
$1\cdot W_1W_0 = \frac\pi2$。

\textbf{14.} $W_{2n+1} \leq W_{2n} \leq W_{2n-1}$（由 $\sin^n$ 的逐点
单调性），且 $\frac{W_{2n-1}}{W_{2n+1}} =
\frac{2n+1}{2n} \to 1$，夹逼得 $\frac{W_{2n}}{W_{2n+1}} \to
1$。结合 $W_{2n}W_{2n+1} = \frac{\pi}{2(2n+1)}$
（第 13 题）：$W_{2n}^2 \sim \frac\pi{4n}$，故 $W_{2n} \sim
\frac12\sqrt{\frac\pi n}$，且 $\binom{2n}n4^{-n} =
\frac2\pi W_{2n} \sim \frac1{\sqrt{\pi n}}$。

\textbf{15.} 第 9--10 题从未用到该常数的值。
令 $e_n = \ln n! - (n+\frac12)\ln n + n$，则
差 $e_n - e_{n+1}$ 落在 $\bigl(0, \frac1{12n} -
\frac1{12(n+1)}\bigr)$ 中，故 $(e_n)$ 递减，而 $(e_n -
\frac1{12n})$ 递增：二者为相邻数列，
收敛于公共极限 $\ell$。于是 $n! \sim K n^{n+1/2}\eu^{-n}$，$K =
\eu^\ell$。

\textbf{16.} 将含未知常数的斯特林公式代入
中心二项式：
\[
\binom{2n}n4^{-n}\sqrt n \sim
\frac{K\,(2n)^{2n+1/2}\eu^{-2n}}{\bigl(K\,n^{n+1/2}
\eu^{-n}\bigr)^2}\,4^{-n}\sqrt n
= \frac{2^{2n}\sqrt{2}\,K\,n^{2n+1/2}}{K^2\,n^{2n+1}}
\,4^{-n}\sqrt n = \frac{\sqrt2}{K},
\]
由第 14 题可得 $\frac{\sqrt2}K = \frac1{\sqrt\pi}$，
即 $K = \sqrt{2\pi}$。第三至四部分由此
从头重证斯特林公式；将其代入第一部分的恒等式，
算出 $\int_\R\eu^{-u^2/2}\dd u = \sqrt{2\pi}$，无需
借助\omterm{ex:b3:product:polar}{极坐标}：沃利斯与高斯相互印证。

\textbf{17.} $\frac{\binom{2n}{n+j}}{\binom{2n}n} =
\frac{(n!)^2}{(n+j)!\,(n-j)!} =
\prod_{i=1}^{j}\frac{n-i+1}{n+i}$ 对 $j \geq 0$ 成立（
由对称性，对 $j < 0$ 亦然）。取对数，其中 $1 \leq i
\leq j \leq K\sqrt n$：
\[
\ln\frac{n-i+1}{n+i} = \ln\Bigl(1 - \frac{i-1}n\Bigr) -
\ln\Bigl(1 + \frac in\Bigr) = -\frac{2i-1}{n} +
O\Bigl(\frac{i^2}{n^2}\Bigr),
\]
又 $\sum_{i\leq j}(2i - 1) = j^2$，而误差之和为
$O(j^3/n^2) = O(n^{-1/2})$，一致地有
$\exp\bigl(-\frac{j^2}n + O(n^{-1/2})\bigr)$.

\textbf{18.} $\eu^{-n}\frac{n^n}{n!} \sim \eu^{-n}
\frac{n^n}{\sqrt{2\pi n}\,n^n\eu^{-n}} =
\frac1{\sqrt{2\pi n}}$。均值为 $n$ 的泊松随机变量
的标准差为 $\sqrt n$，而 $\frac1{\sqrt{2\pi n}}$ 恰为
高斯分布的峰值高度 $\frac1{\sigma\sqrt{2\pi}}$，
即局部中心极限定理，在众数处已有预示。

\textbf{19.} 对 $a \in \intoo01$，由
\cref{pb:b3:lebesgue:1}（即其中第 14 题）的斜率引理，用于凸函数 $\log\Gamma$ 在 $n$ 附近，
可得 $(n-1)^a \leq
\frac{\Gamma(n+a)}{\Gamma(n)} \leq n^a$，其与 $n^a$ 的比值
被 $(1 - \frac1n)^a \to 1$ 夹逼。设 $a = m + a'$
（$m \in \N$，$a' \in \intco01$）：$\Gamma(n+a) = (n + a - 1)
\cdots(n + a')\Gamma(n + a')$，其中 $m$ 个因子中每一个都是
$n(1 + O(\frac1n))$：将上述估计相乘即可。

\textbf{20.} 递推式 $v_d = \frac{2\pi}dv_{d-2}$
（据 $\Gamma(\frac d2 + 1) = \frac d2\Gamma(\frac d2)$ 而来）
从 $v_1 = 2$，$v_2 = \pi$ 出发，可得：
\[
v_3 = \frac{4\pi}3,\quad v_4 = \frac{\pi^2}2,\quad
v_5 = \frac{8\pi^2}{15},\quad v_6 = \frac{\pi^3}6,\quad
v_7 = \frac{16\pi^3}{105}.
\]
比值 $\frac{2\pi}d$ 大于 $1$，恰在 $d \leq 6$ 时成立，
故每种奇偶性下均先增后减，数值上 $v_4
\approx 4.93$，$v_5 \approx 5.26$，$v_6 \approx 5.17$，
整体最大值在 $d = 5$ 处取得。生成函数为 $v_{2k} =
\frac{\pi^k}{k!}$，于是 $\sum_kv_{2k}x^{2k} =
\eu^{\pi x^2}$——将所有偶数维球体的体积
合并为一个指数，并立即给出一个快于几何级数的衰减
估计，适用于 $v_d$。

\textbf{21.} 分子、分母中均用斯特林公式，并取 $k =
\alpha n$：
\[
\binom n{\alpha n} \sim
\frac{\sqrt{2\pi n}\,n^n}
{\sqrt{2\pi\alpha n}\,(\alpha n)^{\alpha n}\,
\sqrt{2\pi(1-\alpha)n}\,((1-\alpha)n)^{(1-\alpha)n}}
= \frac{\eu^{nH(\alpha)}}{\sqrt{2\pi\alpha(1-\alpha)n}},
\]
由于 $n^n/(\alpha n)^{\alpha n}((1-\alpha)n)^{(1-\alpha)n} =
\alpha^{-\alpha n}(1-\alpha)^{-(1-\alpha)n} =
\eu^{nH(\alpha)}$（$n$ 的幂次相消：$\alpha n +
(1-\alpha)n = n$），$\eu^{-n}$ 同样相消。当
$\alpha = \frac12$ 时：$H = \ln2$，前因子为
$\sqrt{2/(\pi n)}$——又见第 5 题。
$H$ 严格凹（其二阶导数为 $-\frac1{\alpha(1-\alpha)} < 0$），
使其最大值 $\ln 2$ 仅在 $\alpha = \frac12$ 处取得。当
$\alpha \neq \frac12$ 时，$\binom n{\alpha n}2^{-n} \approx
\eu^{-n(\ln2 - H(\alpha))}$ 指数衰减——这正是
每一条有关抛硬币的集中性陈述背后的
组合机制。

\textbf{22.} 由 $s_{d-1} = dv_d$ 与第 20 题得：
\[
s_0 = 2,\ \ s_1 = 2\pi,\ \ s_2 = 4\pi,\ \ s_3 = 2\pi^2,\ \
s_4 = \frac{8\pi^2}3,\ \ s_5 = \pi^3,\ \ s_6 =
\frac{16\pi^3}{15},
\]
数值上 $2,\ 6.28,\ 12.57,\ 19.74,\ 26.32,\ 31.01,\
33.07$；且 $s_7 = \frac{\pi^4}3 \approx 32.47 < s_6$，
最大者为 $6$ 维球面。递推式 $s_{d+1} =
(d+2)\,v_{d+2} = (d+2)\,\frac{2\pi}{d+2}\,v_d = 2\pi v_d =
\frac{2\pi}d\,s_{d-1}$ 表明，与体积的情形相同，由 $\frac{2\pi}d$ 驱动的
上升与快于几何级数的下降。维数超过
七之后，球面的收缩速度快于任何几何
数列。

\textbf{23.} 第 11 题恰给出 $\frac1{12n+1} < d_n <
\frac1{12n}$，且
\[
\frac1{12n} - \frac1{12n+1} = \frac1{12n(12n+1)} =
O\Bigl(\frac1{n^2}\Bigr),
\]
于是 $d_n = \frac1{12n} + O(\frac1{n^2})$，且 $\eu^{d_n} = 1 +
\frac1{12n} + O(\frac1{n^2})$；
乘以 $\sqrt{2\pi n}(n/\eu)^n$ 即得修正后的公式。当 $n =
10$：$\sqrt{20\pi}\,(10/\eu)^{10} = 7.92665 \times 453999.3
\approx 3\,598\,696$，偏低 $30\,104$（相对误差
$8.3\cdot10^{-3}$）；乘以 $1 + \frac1{120}$ 得
$3\,628\,685$，偏低 $115$（相对误差
$3.2\cdot10^{-5}$）。夹逼本身将 $10!$ 限定在
$3\,598\,696\,\eu^{1/121} \approx 3\,628\,559$ 与
$3\,598\,696\,\eu^{1/120} \approx 3\,628\,808$——上界
在前七位数字上仅偏差八个单位。

\textbf{24.} 第一部分的换元 $x = t + \sqrt t\,u$，
作用于截断后的积分，得
\[
\int_0^{t} x^{t}\eu^{-x}\,\dd x
= t^{t}\eu^{-t}\sqrt t\int_{-\sqrt t}^{0} g_t(u)\,\dd u ,
\]
而范围 $0 \leq x \leq t$ 化为 $-\sqrt t \leq u \leq
0$。第 3 题中的控制函数也覆盖 $g_t\mathbf 1_{u < 0}$，
故由控制收敛定理得
\[
\int_{-\sqrt t}^{0}g_t(u)\,\dd u \longrightarrow
\int_{-\infty}^{0}\eu^{-u^2/2}\dd u = \frac{\sqrt{2\pi}}2 ,
\]
而第 4 题给出 $\Gamma(t+1) \sim
t^t\eu^{-t}\sqrt t\,\sqrt{2\pi}$。比值趋于
$\frac12$。从概率角度看，形状参数较大的伽马随机变量
渐近地将其一半质量放在其众数的每一侧
——中心极限定理的对称性，
由一次换元便可读出。

\textbf{25.} 设 $\lambda = \frac{\alpha}{1-\alpha} \in
\intoc01$。对于 $k \leq \lfloor\alpha n\rfloor \leq \alpha n$，
有 $\lambda^{k - \alpha n} \geq 1$，因此
\[
\sum_{k=0}^{\lfloor\alpha n\rfloor}\binom nk
\leq \lambda^{-\alpha n}\sum_{k=0}^{n}\binom nk\lambda^{k}
= \Bigl(\lambda^{-\alpha}(1 + \lambda)\Bigr)^{n},
\]
且由 $\lambda = \frac\alpha{1-\alpha}$ 可得
\[
\lambda^{-\alpha}(1+\lambda)
= \alpha^{-\alpha}(1-\alpha)^{\alpha}\cdot\frac1{1-\alpha}
= \alpha^{-\alpha}(1-\alpha)^{-(1-\alpha)}
= \eu^{H(\alpha)} .
\]
最优性：极小化 $f(\lambda) = -\alpha\ln\lambda +
\ln(1+\lambda)$，在 $\lambda > 0$ 上，方程 $f'(\lambda)
= -\frac\alpha\lambda + \frac1{1+\lambda} = 0$ 的唯一解为
$\lambda = \frac\alpha{1-\alpha}$，且由于
$f'' > 0$，它是极小值点——即指数倾斜（切尔诺夫）的选取。
衔接：由第 21 题，单项 $k =
\lfloor\alpha n\rfloor$ 的项已达到
$\eu^{nH(\alpha)}/\sqrt{2\pi\alpha(1-\alpha)n}$ 的量级，因此
\[
\frac{\eu^{nH(\alpha)}}{C\sqrt n} \leq
\sum_{k\leq\alpha n}\binom nk \leq \eu^{nH(\alpha)} :
\]
速率 $H(\alpha)$ 是精确的，整个和式至多比其最大项多出
一个 $\sqrt n$ 因子。除以 $2^n$ 后，
这就是公平硬币的尾部界 $\P(S_n \leq \alpha n) \leq
\eu^{-n(\ln2 - H(\alpha))}$——\omterm{def:b3:measure:measure}{测度}集中现象，
只需一行即可读出。
\end{solution}
